\section{Known limits, stated before the data: full list} \label{app:limits}

\begin{itemize}
\item The archival sectors were inspected, and one choice in the ephemeris (B modelled jointly with A,
  Section~\ref{sec:predictions}) was made after looking at them.
\item With the injection-calibrated errors, the archival transit times scatter less than expected ($\chi^2 = 8.3$ for
  20 degrees of freedom). The prediction uncertainty of A may therefore be somewhat too large, which makes the timing
  rule more lenient.
\item $D_{98}$ for A (450\,ppm), the archival centre from which $\dhi$ is built, is close to the primary expectation
  (442\,ppm). For A the consistency band is therefore set mostly by the noise and by the two expectation errors, not by
  an excess in the discovery sector.
\item The edge-zone rule was introduced during the adversarial audit of a first draft, after the archival dry runs had
  shown the artefacts at data gaps. Before the rule was fixed, two zone widths (6\,h and 9\,h; at that stage measured from the
  nearest data gap) were tried once on archival data to see whether the artefacts were the cause. Both gave the same
  result (Sector 32: $\ddet = 211$ and 214\,ppm against 236 without zones; Sector 98: $\ddet = 142$ and 141 against
  581; no real transit lost), so the trial could not favour one width. The adopted width is $W/2$ by the definition of
  the filter, and it scales with $W$ in the variants. The archival recovery of A and B (Appendix~\ref{app:validation}) is
  a validation made after the rule was fixed, not the reason for it.
\item The change of percentiles (Appendix~\ref{app:history}, row 3b) was also made after archival measurements, to bring
  the measured false-alarm level down. It makes detection stricter and widens the consistency band slightly. The percentiles were set once
  (99.9 and 0.1) and then measured; they were not tuned further. The measured levels (SPOC 2-min: A 0.47\,\%, B
  0.49\,\% overall, with the thresholds of the Sector 110 run since audit 6; Appendix~\ref{app:fa}) pool every trial set
  that was run, including audit 3's. An intermediate version of this document quoted a lower level for A from one set
  only and a grouping of $n$ chosen after seeing the counts; both were corrected after audit 4.
\item The equal-stack localisation scheme was checked on audit 4's synthetic Sector 110 pixel file (a 2-d gap and 8
  transits imposed). The on-target cloud has standard deviations of 3.7 to $3.9''$ (East) and 4.7 to $5.2''$ (North),
  26 to 27 phase slots and no failed trial, against 3.4 to $3.5''$ and 9.2 to $10.7''$ with the earlier
  phase-rejection scheme. This check was made outside this deposit on a synthetic file built from the Sector 32 pixels;
  its scripts and outputs are kept by me and are not part of this deposit, so it is given for information only.
\item Edge zones cost data: in a Sector 110 light curve with two or three segments, a transit has roughly a 6 to 10\,\%
  chance to fall in an edge zone and not be tested.
\item Neighbouring windows of the noise pool overlap (steps of $\Tfourteen/2$), so the pool has fewer independent
  windows than its size. This does not bias the noise distribution, but its far tail rests on few independent windows.
  This is why the false-detection levels are measured separately (Appendix~\ref{app:fa}).
\item The archival localisation inputs include Sector 98 injection trials that overlap signal B (113 of 300). Using only
  the clean trials would change the archival combined $p$ from 0.061 to 0.053. The headline values are used as
  published. They come from the item 11 recipe, whose fakes may sit next to data gaps; the Sector 110 clouds exclude
  such fakes and are narrower. For the same Sector 32 dry-run offset $(+5.46, +0.38)''$ the Gaussian $p$ was 0.29 with
  the first wrapper (unequal stacks, fakes anywhere), 0.23 with equal stacks but no edge-zone rule, and 0.093 with the
  final scheme; the empirical $p$ of the final rule is 0.156. The combined value mixes the two kinds of cloud and uses
  Gaussian tails, and it is descriptive only.
\item The localisation $p$ was changed after audit 5 from the Gaussian tail $\exp(-\delta^2/2)$ to the empirical $p$,
  because archival on-target injections showed that the Gaussian tail is too light (Sector 32: 12 of 300 below
  0.0027). The thresholds 0.05 and 0.0027 were not changed, and the method was fixed before its value for the Sector 32
  dry run was computed. For that offset it raises $p$ from 0.093 to 0.156: such offsets are less unusual among
  on-target injections than a Gaussian suggests. For moderate offsets the final rule is therefore more lenient than the
  Gaussian one, while a false ``off target'' for an on-target signal now occurs near the nominal rate
  (Appendix~\ref{app:validation}). On audit 4's synthetic Sector 110 pixel file (8 transits, a 2-d gap; a check made
  outside this deposit, kept by me, for information only) the same coverage check gives 4.5 and 5.6\,\% below 0.05 and 0.07 and 0.50\,\% below
  0.0027 (Gaussian tail: 7.4 to 7.6\,\% and 1.7 to 1.9\,\%). All injections of one sector share pixels and overlapping
  windows, so the empirical $p$ assumes that a real on-target signal behaves like the injections at random phases. Its
  resolution is $1/(N_{\rm cal}+1)$, or 1/3001 when all 3000 calibration trials succeed. Near 0.0027 only about 8 calibration trials lie further out, so its Monte Carlo error there is
  about a third of its value.
\item Box depths are lower than fitted transit depths (limb darkening, finite exposure). The expected and measured
  values here are both box depths, so they can be compared with each other, but not with fitted depths quoted
  elsewhere.
\end{itemize}

